\documentclass[11pt]{article}
\usepackage{amsmath,amssymb}
\usepackage[hidelinks]{hyperref}
\title{Clearance and connectivity after a corridor crash}
\date{24 November 2026}
\begin{document}
\maketitle

\paragraph{Question.} After a robot loses mobility and radio in a single-file
corridor, can the surviving radio chain reconnect by time \(T\) without
crossing or approaching the stopped body too closely? The answer below
is exact for one known fault position, fixed order and endpoints, and
single-integrator survivors with a common speed bound \(v>0\).
There is no claim about dynamics with acceleration limits, equipment in
the field, lawfully obtaining the initial configuration, or the full R6
obligation.

\paragraph{Source model.} Work on the line segment \([0,L]\).
The \(n\geq 2\) survivors start at ordered positions \(x_0,\dots,x_{n-1}\),
with \(x_0=0\), \(x_{n-1}=L\). A disabled robot remains at \(f\in(0,L)\),
has no active radio, and must be separated from every survivor by \(g>0\).
Adjacent survivors must also be separated by at least \(g\); these lower
bounds must hold throughout. The initial physical configuration satisfies
these bounds. Radios on surviving units have range \(r\), with no
directionality or occlusion; the stopped body does not block radio
propagation. A link exists exactly when the distance is at most \(r\).
A sorted surviving chain is connected exactly when every adjacent gap is
at most \(r\): if an adjacent gap exceeds \(r\), all links across that cut
also exceed \(r\). Connectivity may be lost at the fault; the requirements
here are restoration by \(T\) and maintenance afterward by holding still,
not continuous post-fault connectivity before restoration.

\paragraph{Exact recovery test.} If \(r<g\), no connected safe terminal configuration exists.
Otherwise let \(D_i=[\ell_i,u_i]\) be the intersection of the segment,
\([x_i-vT,x_i+vT]\), the original side's safety half-line
\((-\infty,f-g]\) or \([f+g,\infty)\), and the fixed-position singleton at
an endpoint. Empty intervals are allowed as failure indicators.
Set \(J_0=D_0\). For each later vertex let
\[
  J_i = D_i\cap [\inf J_{i-1}+g,\ \sup J_{i-1}+r],
\]
propagating emptiness, never taking endpoints of an empty interval.
Then safe recovery by \(T\) exists if and only if \(J_{n-1}\ne\emptyset\).
The scan costs \(O(n)\) arithmetic operations, and deciding existence
needs only the current interval. Reconstruction retains the prefix
intervals and costs \(O(n)\) further operations.

For the induction, \(J_{i-1}\) is exactly the set of achievable endpoints
for the preceding chain. For nonempty \(J_{i-1}=[a,b]\), a candidate \(y\)
can extend the chain exactly when
\([a,b]\cap[y-r,y-g]\) is nonempty. This is equivalent to
\(a+g\leq y\leq b+r\), since \(a\leq b\) and \(g\leq r\).
Choose \(y_{n-1}\in J_{n-1}\) and recursively choose
\(y_{i-1}=\max(\inf J_{i-1},y_i-r)\). The same argument gives exact
integer arithmetic when all bounds and candidate positions are integral.

Every physically admissible motion satisfies the displacement and
terminal inequalities. Conversely, for \(T>0\), take the straight
interpolation \(x_i(t)=(1-t/T)x_i+(t/T)y_i\). 
The displacement bound gives speed at most \(v\); every clearance,
side and endpoint constraint is affine and holds at both ends, hence
throughout. For \(T=0\) the test checks the initial state directly.
The terminal upper gaps are not inferred at intermediate times from
one connected endpoint. If the initial upper gaps also hold, they are
preserved by the same affine argument. Thus there is neither a hidden
violation between sample instants nor an assumption that a fault
already violating connectivity can be undone retroactively.

\paragraph{Matched cases.} With survivors at \(0,2,7,9\), stopped body
at \(4\), \(g=1\), and \(v=1\), range \(3\) yields the tight deadline
\(T=1\): the three gaps must all equal \(3\), requiring the terminal
positions \(0,3,6,9\). Increasing range to \(4\) yields \(T=1/2\):
the initially five-unit middle gap can shrink by at most \(2T\), and
\(0,5/2,13/2,9\) attains the bound. Separately, if there are survivors on
both sides and \(r<2g\), the stopped body makes a radio connection
physically impossible: the cross-body gap is at least \(2g\).
For \(L=6\), \(g=1\), \(r=3/2\), the initially connected formation
\(0,1,2,3,9/2,6\) therefore cannot reconnect across a crash at \(3\).
If the failed body is removed instead, the same survivors can occupy
\(0,3/2,3,9/2,6\). This compares physical continuation classes, not
policies with different access to information.

\paragraph{Evidence and scope.} The direct comparison solves the same
terminal constraints as a linear program and compares recovered
deadlines and rational certificates in 480 cases. Integer exhaustive
enumeration agrees in 8,928 cases. \texttt{CorridorFormation.lean}
checks exact integer endpoint propagation for all chain lengths and
the real affine interpolation and speed lemmas. Lean does not check
the Python code, the radio correspondence, or a distributed sensing
protocol. This is a reachability oracle for a coordination layer that
can lawfully obtain its inputs. Standard chain constraint elimination
can implement the same scan; no historical novelty or gain over a
well-factored direct chain solver is claimed.

Maeda, Endo and Matsuno give
\href{https://doi.org/10.1016/j.jfranklin.2023.05.013}{decentralized
crash recovery and avoidance of the disabled robot}; their local
controller is not replayed here. The bounded single-file oracle is a
complementary exact reference, not a claim that it replaces their
controller. \href{https://arxiv.org/abs/2609.29939}{Managed robot
departure} can prepare replacement links before removal, unlike the
unannounced loss in this model.
\href{https://arxiv.org/abs/2210.01354}{Barrier-function composition}
addresses control viability under simultaneous constraints; our
linear interpolation result relies on the stated affine convexity.
This reference does not discharge R6 or its connection to the shared
scientific laws. Hidden crash position, radio occlusion, acceleration,
branching corridors and evidence to distinguish physical infeasibility
from local uncertainty remain separate checks.
\end{document}
